\documentclass[10pt]{article}

\usepackage[margin=0.75in]{geometry}
\usepackage{amsmath,amsthm,amssymb,bm}
\usepackage{xcolor}
\usepackage{cancel}
\usepackage{graphicx}
\usepackage{changepage}
\usepackage{circuitikz}
\usepackage{pgfplots}
\usepackage{physics}
\usepackage{hyperref}
\usepackage{siunitx}
\usepackage{minted}
\usepackage{fontspec}
\usepackage{relsize}
\usepackage{subfig}
\usepackage{mhchem}
\usepackage{todonotes}
\usepackage{biblatex}
\usepackage{multicol, multirow, booktabs}
\usepackage[breakable]{tcolorbox}
\usepackage[inline]{enumitem}
\patchcmd{\thebibliography}{\section*{\refname}}{}{}{}
\addbibresource{sources.bib}

\theoremstyle{definition}
\newtheorem{problem}{Problem}
\newtheorem{soln}{Solution}

\pgfplotsset{compat=newest}
\usetikzlibrary{lindenmayersystems}
\usetikzlibrary{arrows}
\usetikzlibrary{calc}
\usetikzlibrary{positioning, fit}
\usetikzlibrary{3d, perspective}

\definecolor{incolor}{HTML}{303F9F}
\definecolor{outcolor}{HTML}{D84315}
\definecolor{cellborder}{HTML}{CFCFCF}
\definecolor{cellbackground}{HTML}{F7F7F7}
\newcommand{\ui}{\hat{i}}
\newcommand{\uj}{\hat{j}}
\newcommand{\uk}{\hat{k}}
\newcommand{\ux}{\hat{x}}
\newcommand{\uy}{\hat{y}}
\newcommand{\uz}{\hat{z}}
\newcommand{\uv}[1]{\hat{\bm{{#1}}}}
\newcommand{\pr}[1]{{#1^\prime}}
\newcommand{\justif}[2]{&{#1}&\text{#2}}
\newcommand{\dul}[1]{\underline{\underline{#1}}}
\newcommand{\pdiff}[2][]{{\frac{\partial^{#1}}{\partial {{#2}^{#1}}}}}
% \newcommand{\dif}[2][]{{\frac{d^{#1}}{d{{#2}^{#1}}}}}
\pgfdeclarelayer{background}  
\pgfsetlayers{background,main}
\AtBeginDocument{\RenewCommandCopy\qty\SI}

\makeatletter
\newcommand{\boxspacing}{\kern\kvtcb@left@rule\kern\kvtcb@boxsep}
\makeatother
\newcommand{\prompt}[4]{
    \ttfamily\llap{{\color{#2}[#3]:\hspace{3pt}#4}}\vspace{-\baselineskip}
}

\newcommand{\thevenin}[2]{
  \begin{center}
    \begin{circuitikz} \draw
      (0,0) -- (2,0) to[battery1, l_=$V_{Th}\eq#1$] (2,2) 
      to[resistor, l_=$R_{Th}\eq#2$] (0,2)
      ;
      \draw [o-] (-.07,2.079);
      \draw [o-] (-.07,0.079);
    \end{circuitikz}
  \end{center}
}

\newcommand{\norton}[2]{
  \begin{center}
    \begin{circuitikz} \draw
      (0,0) -- (3,0) to[american current source, l_=$I_{N}\eq#1$] (3,2) -- (0,2) (2,0)
      to[resistor, l=$R_{N}\eq#2$] (2,2)
      ;
      \draw [o-] (-.07,2.079);
      \draw [o-] (-.07,0.079);
    \end{circuitikz}
  \end{center}
}

\newcommand{\highlight}[1]{\colorbox{yellow}{#1}}

\newcommand{\ti}[1]{\widetilde{#1}}

\newfontface{\Kaufmann}{Kaufmann}
\DeclareTextFontCommand{\kf}{\Kaufmann}
\newcommand{\scriptr}{\fontsize{12pt}{12pt}\kf{r}}

\newfontface{\KaufmannB}{Kaufmann Bold}
\DeclareTextFontCommand{\kfb}{\KaufmannB}
\newcommand{\bscriptr}{\fontsize{12pt}{12pt}\kfb{r}}

\newcommand{\bv}[1]{\vec{#1}}

\title{Physics 4605H: Assignment VI}
\author{Jeremy Favro (0805980) \\ Trent University, Peterborough, ON, Canada}
\date{\today}

\begin{document}
\maketitle

% PROBLEM 1
\begin{problem}
Consider the case of a nuclear spin in the presence of an applied field
$$\bv{B}(t)=U_0\uz+B_{rf}\cos\left(\omega_{rf}t-\phi\right)\ux$$
In the lab frame,
$$\hat{H}=-\hat{\bv{\mu}}\cdot\bv{B}=-\frac{\hbar \omega_L}{2}\hat{\sigma}_z-\hbar\omega_R\cos\left(\omega_{rf}t-\phi\right)\hat{\sigma}_x.$$
Your goal in this problem is to find the Hamiltonian in the rotating reference frame $\hat{H}'$ in
this case.
Recall that to switch to the rotating frame we did a unitary transformation using the operator
$$\hat{R}_z(\omega_{rf}t)=e^{-i\omega_{rf}t\hat{\sigma}_z/2}\text{ such that }\ket{\psi'(t)}=\hat{R}_z(\omega_{rf}t)\ket{\psi(t)}$$
(\highlight{note that I've added a $t$ to the exponent, this was not present in the assignment}).
Then, we took the time derivative of $\ket{\psi'(t)}$
$$i\hbar\frac{\partial}{\partial t}\ket{\psi'(t)}=\hat{H}'\ket{\psi'(t)}$$
and we found
$$\hat{H}'=\left\{i\hbar\left(\frac{\partial}{\partial t}\hat{R}_z(\omega_{rf}t)\right)+\hat{R}_z(\omega_{rf}t)\hat{H}\right\}\hat{R}^\dagger_z(\omega_{rf}t)$$

Your job is to evaluate this for the Hamiltonian $\hat{H}$ above. In particular you should find that
in the rotating wave approximation and when $\omega_{rf}=\omega_{L}$,
$$\hat{H}'=-\frac{\hbar\omega_R}{2}\left(\cos\phi\hat{\sigma}_x+\sin\phi \hat{\sigma}_y\right).$$
You may use (without proof) the following results which we derived in class:
\begin{align*}
  \hat{R}_z(\omega_{rf}t)\hat{\sigma}_z\hat{R}_z^\dagger(\omega_{rf}t) & =\hat{\sigma}_z                                                    \\
  \hat{R}_z(\omega_{rf}t)\hat{\sigma}_x\hat{R}_z^\dagger(\omega_{rf}t) & =\cos(\omega_{rf}t)\hat{\sigma}_x+\sin(\omega_{rf}t)\hat{\sigma}_y
\end{align*}
You will also want to use expressions for cosine and sine of the sum of two angles and related
trigonometric identities.
\end{problem}
\begin{soln}
  To make this a bit easier to write out we break down the given expression for $\hat{H}'$,
  $$\hat{H}'=\left\{i\hbar\left(\frac{\partial}{\partial t}\hat{R}_z(\omega_{rf}t)\right)+\hat{R}_z(\omega_{rf}t)\hat{H}\right\}\hat{R}^\dagger_z(\omega_{rf}t)$$
  as two terms. The first is
  $$i\hbar\left(\frac{\partial}{\partial t}\hat{R}_z(\omega_{rf}t)\right)\hat{R}^\dagger_z(\omega_{rf}t)$$
  which expands out to look like
  $$i\hbar\left[\left(\frac{-i\omega_{rf}\hat{\sigma}_z}{2}\right)\left(e^{-i\omega_{rf}t\hat{\sigma}_z/2}\right)e^{i\omega_{rf}t\hat{\sigma}_z/2}\right]
    =\frac{\hbar\omega_{rf}}{2}\hat{\sigma}_z.
  $$
  The second term expands out to
  $$
    \hat{R}_z(\omega_{rf}t)\hat{H}\hat{R}^\dagger_z(\omega_{rf}t)
    =e^{-i\omega_{rf}t\hat{\sigma}_z/2}\left[-\frac{\hbar \omega_L}{2}\hat{\sigma}_z-\hbar\omega_R\cos\left(\omega_{rf}t-\phi\right)\hat{\sigma}_x\right]e^{i\omega_{rf}t\hat{\sigma}_z/2}
  $$
  and bringing the $\hat{R}_z$ operators inside the brackets,
  \begin{align*}
     & =-\frac{\hbar \omega_L}{2}\hat{R}_z(\omega_{rf}t)\hat{\sigma}_z\hat{R}^\dagger_z(\omega_{rf}t)-\hat{R}_z(\omega_{rf}t)\hbar\omega_R\cos\left(\omega_{rf}t-\phi\right)\hat{\sigma}_x\hat{R}^\dagger_z(\omega_{rf}t)   \\
     & =-\frac{\hbar \omega_L}{2}\hat{\sigma}_z-\hbar\omega_R\cos\left(\omega_{rf}t-\phi\right)\left[\cos(\omega_{rf}t)\hat{\sigma}_x+\sin(\omega_{rf}t)\hat{\sigma}_y\right]\justif{\quad}{Applying the given identities}.
  \end{align*}
  Now recombining the first and second terms we get that
  \begin{align*}
    \hat{H}' & =\frac{\hbar\omega_{rf}}{2}\hat{\sigma}_z
    -\frac{\hbar \omega_L}{2}\hat{\sigma}_z-\hbar\omega_R\cos\left(\omega_{rf}t-\phi\right)\left[\cos(\omega_{rf}t)\hat{\sigma}_x+\sin(\omega_{rf}t)\hat{\sigma}_y\right] \\
             & =\frac{\hbar\left(\omega_{rf}-\omega_L\right)}{2}\hat{\sigma}_z
    -\hbar\omega_R\cos\left(\omega_{rf}t-\phi\right)\left[\cos(\omega_{rf}t)\hat{\sigma}_x+\sin(\omega_{rf}t)\hat{\sigma}_y\right]
  \end{align*}
  which in the limit of the first approximation, $\omega_{rf}=\omega_L$ becomes
  $$
    \hat{H}'=-\hbar\omega_R\cos\left(\omega_{rf}t-\phi\right)\left[\cos(\omega_{rf}t)\hat{\sigma}_x+\sin(\omega_{rf}t)\hat{\sigma}_y\right]
  $$

  Now we use the $\cos(a-b)=\cos a\cos b+\sin a\sin b$ identity,
  \begin{align*}
     & =-\hbar\omega_R\left[\cos(\omega_{rf}t)\cos(\phi)+\sin(\omega_{rf}t)\sin(\phi)\right]\left[\cos(\omega_{rf}t)\hat{\sigma}_x+\sin(\omega_{rf}t)\hat{\sigma}_y\right] \\
     & =-\hbar\omega_R\left[\left[\cos^2(\omega_{rf}t)\cos(\phi)+\cos(\omega_{rf}t)\sin(\omega_{rf}t)\sin(\phi)\right]\hat{\sigma}_x
      +\left[\sin(\omega_{rf}t)\cos(\omega_{rf}t)\cos(\phi)+\sin^2(\omega_{rf}t)\sin(\phi)\right]\hat{\sigma}_y\right]
  \end{align*}
  and we use the $\cos^2(x)=\frac{1}{2}(1+\cos(2x))$ and $\sin^2(x)=\frac{1}{2}(1-\cos(2x))$ identities,
  \begin{align*}
     & =-\hbar\omega_R\left[\left[\frac{1}{2}(1+\cos(2\omega_{rf}t))\cos(\phi)+\cos(\omega_{rf}t)\sin(\omega_{rf}t)\sin(\phi)\right]\hat{\sigma}_x\right. \\
     & \quad\left.+\left[\sin(\omega_{rf}t)\cos(\omega_{rf}t)\cos(\phi)+\frac{1}{2}(1-\cos(2\omega_{rf}t))\sin(\phi)\right]\hat{\sigma}_y\right]
  \end{align*}
  and we use $\sin(x)\cos(x)=\frac{1}{2}\sin(2x)$,
  \begin{align*}
     & =-\frac{\hbar\omega_R}{2}\left[\left[(1+\cos(2\omega_{rf}t))\cos(\phi)+\sin(2\omega_{rf}t)\sin(\phi)\right]\hat{\sigma}_x
    +\left[\sin(2\omega_{rf}t)\cos(\phi)+(1-\cos(2\omega_{rf}t))\sin(\phi)\right]\hat{\sigma}_y\right]                           \\
     & =-\frac{\hbar\omega_R}{2}\left[\left[\cos(\phi)+\cos(2\omega_{rf}t)+\sin(2\omega_{rf}t)\sin(\phi)\right]\hat{\sigma}_x
    +\left[\sin(2\omega_{rf}t)\cos(\phi)+\sin(\phi)-\cos(2\omega_{rf}t)\right]\hat{\sigma}_y\right]                          \\
    & =-\frac{\hbar\omega_R}{2}\left[\cos(\phi)\hat{\sigma}_x+\sin(\phi)\hat{\sigma}_y+\left[\cos(2\omega_{rf}t)+\sin(2\omega_{rf}t)\sin(\phi)\right]\hat{\sigma}_x
    +\left[\sin(2\omega_{rf}t)\cos(\phi)-\cos(2\omega_{rf}t)\right]\hat{\sigma}_y\right]                          
  \end{align*}
  we can now make the actual rotating-wave approximation. I like the way that the book explains it using the fact that the contribution of the 
  $2\omega_{rf}t$ terms to the time evolution operator ($\exp(-i\int\hat{H}'dt)$) is bounded for $t\gg 1/\omega_{rf}$ but the contribution of the static (in time) phase terms increases linearly with time,
  eventually dominating over the effect of the rotating terms. Hence we can neglect the overall contribution of the higher frequency terms which leads to the desired result,
  $$\hat{H}' =-\frac{\hbar\omega_R}{2}\left[\cos(\phi)\hat{\sigma}_x+\sin(\phi)\hat{\sigma}_y\right].$$
\end{soln}

\end{document}